% source: manti.01.ma, 1ère SM, série fonctions (Drive 0B6pfSJcEVng0bHVtbDUzcTB4N1k)
\ex{Fonctions}{bHVt}{Ex. 1}{5}
Déterminer l'ensemble de définition de la fonction $f$ dans chacun des cas suivants :
\begin{enumerate*}[label=\arabic*),itemjoin={\quad}]
\item $f(x)=\dfrac{x}{2|x+2|-1}$
\item $f(x)=\dfrac{x+1}{4x^2-1}$
\item $f(x)=\dfrac{2x-1}{|2x+3|+x}$
\item $f(x)=\dfrac{3x-2}{4\sqrt x-1}$
\item $f(x)=\dfrac{x-1}{3x^2-5x-2}$
\item $f(x)=\dfrac{x}{x|x|-1}$
\item $f(x)=\sqrt{x^2-2x}$
\item $f(x)=\sqrt{x-\dfrac1x}$
\item $\begin{cases}f(x)=\dfrac{x+1}{x^2-2x}&;\ x\geq1\\ f(x)=\dfrac{x^2+x}{3x-2}&;\ x<1\end{cases}$
\end{enumerate*}
\fin
\ex{Fonctions}{bHVt}{Ex. 2}{5}
Étudier la parité de la fonction $f$ dans chacun des cas suivants :
\begin{enumerate*}[label=\arabic*),itemjoin={\quad}]
\item $f(x)=\dfrac{x}{x^2-3|x|}$
\item $f(x)=x+|3x-5|-|3x+5|$
\item $f(x)=\sqrt{x^2-x}+\sqrt{x^2+x}$
\item $f(x)=\dfrac{x^3}{|x+2|+|x-2|+1}$
\item $f(x)=\dfrac{x^3+x^2-2}{x^2-x-2}-\dfrac{x^3-x^2+2}{x^2+x-2}$
\end{enumerate*}
\fin
\ex{Fonctions}{bHVt}{Ex. 3}{8}
\begin{enumerate}
\item Soit $f(x)=\dfrac{x}{x-\sqrt x+1}$. Déterminer l'ensemble de définition de $f$ et montrer que $f$ est majorée par $\dfrac43$ et minorée par $0$.
\item Soit $f(x)=\dfrac{x+2\sin x}{|x|+4}$. Montrer que $f$ est bornée.
\item Montrer que la fonction $f(x)=\dfrac{x^2}{x^2+x+1}$ admet un extremum en $x_0=-2$ en précisant sa nature.
\item Montrer que la fonction $f(x)=\dfrac{x^2}{x^4-x^2+1}$ admet un extremum en $x_0=\sqrt2$ en précisant sa nature \ ([?] point à vérifier).
\item Montrer que la fonction $f(x)=\dfrac{x^3}{x^2-1}$ admet un extremum local en $x_0=\sqrt3$ en précisant sa nature.
\end{enumerate}
\fin
\ex{Fonctions}{bHVt}{Ex. 4}{9}
On considère la fonction $f$ définie par $f(x)=\dfrac{x^2+2}{2x-1}$.
\begin{enumerate}
\item Déterminer l'ensemble de définition de $f$.
\item Montrer que $f$ admet un extremum en $a=2$ en précisant sa nature.
\item Montrer que le taux de variation de $f$ s'écrit $t_f(x,y)=\dfrac{2xy-x-y-4}{(2x-1)(2y-1)}$.
\item Étudier la monotonie de $f$ sur $[2,+\infty[$ et sur $\left]\dfrac12,2\right]$.
\end{enumerate}
\fin
\ex{Fonctions}{bHVt}{Ex. 5}{9}
On considère la fonction $f$ définie par $f(x)=x^3-3x^2-9x$.
\begin{enumerate}
\item Montrer que $f$ admet un extremum en $x_0=-1$ en précisant sa nature.
\item Montrer que $T_f(x,y)=x(x-3)+y(y-3)+xy-9$ et étudier la monotonie de $f$ sur chacun des intervalles $]-\infty,-1]$ et $[3,+\infty[$.
\item Montrer que la fonction $f$ est décroissante sur l'intervalle $[-1,3]$.
\end{enumerate}
\fin
